\subsection{Merge sort}

Merge sort resuelve el problema de ordenamiento (definido en la sección \ref{subsec:problema}).

La idea es dividir el arreglo a ordenar en dos mitades, ordenar cada mitad de forma recursiva (es decir, dividiendo cada una otra vez a la mitad, y así sucesivamente, hasta que queden pedazos de un elemento, que ya están ordenados) y luego combinar las soluciones.

Haciendo explícitas las etapas:
\begin{itemize}
    \item \textbf{Dividir}: Dividir el arreglo en dos mitades.
    \item \textbf{Conquistar}: Ordenar cada mitad usando merge sort (o sea, aplicando estas etapas a cada mitad)
    \item \textbf{Combinar}: Combinar las soluciones de cada mitad, que son subarreglos ordenados.
\end{itemize}

\subsubsection{Dividir y conquistar el arreglo}

Dividir el arreglo no tiene ciencia: se calcula la mitad con división entera y con eso se conocen los índices de cada segmento. La etapa de conquistar es simplemente hacer el llamado recursivo a cada mitad, porque en algún punto al dividir continuamente por la mitad se obtienen arreglos de longitud uno, que ya están ordenados.

Por ahora, para no pensar en cómo combinar las soluciones, encapsulamos ese paso en una función que se implementará después. Con eso, este es el seudocódigo:
\begin{pseudocode}
función merge_sort(números: arreglo, ind_min: entero, ind_max: entero)
    |\br{gi}|if ind_min >= ind_max 
        return|\br{gf}|

    ind_mitad = (ind_min + ind_max) // 2
    merge_sort(números, ind_min, ind_mitad)
    merge_sort(números, ind_mitad+1, ind_max)

    combinar(números, ind_min, ind_mitad, ind_max)
\end{pseudocode}
\begin{annotations}
    \codebrace[xshift=2.6cm, width=7cm]{gi}{gf}{No se retorna nada porque el algoritmo es in-place; se usa el \texttt{return} para detener la función en el caso base de la recursión}
\end{annotations}

Necesitamos tener los índices como parámetros porque indican qué segmento del arreglo se está procesando. La invocación inicial, que resuelve el problema completo, usa los índices extremos del arreglo:

\inline{merge_sort(números, 0, números.longitud-1)}

Nótese que el arreglo no se divide físicamente en cada llamado recursivo; únicamente se delimitan distintas regiones mediante índices. Esto es para que el algoritmo sea \emph{in-place}, es decir, modifique directamente el arreglo original, sin crear y retornar uno nuevo. En cada llamado recursivo ocurre lo siguiente:
\begin{enumerate}
    \item Si el segmento tiene longitud 1 (o 0), se alcanza el caso base y la función termina.
    
    \item Si no, el segmento se divide conceptualmente en dos mitades y se realizan llamados recursivos para la mitad izquierda y luego la derecha. El llamado a la mitad derecha empieza después de que termina el llamado a la mitad izquierda.
    
    \item Una vez ambas llamadas retornan, se ejecuta \inline{combinar}, que mezcla ambas mitades ordenadas in-place dentro del rango delimitado por \inline{ind_min} e \inline{ind_max}.
\end{enumerate}

Para arreglos largos, el caso base solo se alcanza después de muchos llamados recursivos, cuando el arreglo se ha dividido lo suficiente. Es decir: en las primeras invocaciones, siempre ocurre que el segmento se divide conceptualmente. Se acumulan llamados a \inline{merge_sort} uno encima de otro, sin retornar ni combinar. Ese es el \emph{descenso} de la recursión.

Cuando algún llamado alcanza el caso base, ocurre el primer retorno. La ejecución continúa justo después de ese llamado recursivo (o sea, al final de la línea 5 o de la línea 6) y entonces pasa una de dos cosas: o se realiza otro llamado recursivo (que pronto alcanzará el caso base, porque el segmento ya es muy pequeño), o se ejecuta por primera vez el paso de combinar, mezclando esos dos segmentos (que son pequeños, porque hace poco se llegó al caso base). Después de combinar, la función termina y retorna a quien la había invocado. Así, los llamados a \inline{merge_sort} que se habían acumulado empiezan a retornar uno tras otro y la recursión comienza a \emph{desenrollarse}. Ese es el \emph{ascenso} de la recursión.

\subsubsection{Combinar las mitades}

Ahora, solo falta pensar en cómo combinar las mitades. Sorprendentemente, la mejor forma de combinarlas es la forma más intuitiva. Si cada arreglo fuera una pila de cartas, el método natural sería comparar la primera de cada pila y tomar la menor; luego comparar la siguiente de esa pila con la primera de la otra y tomar la menor (que en la combinación es la que le sigue a la primera que se había tomado), y así sucesivamente. El seudocódigo es el siguiente:
\begin{pseudocode}
función combinar(números: arreglo, ind_min: entero, ind_mitad: entero, ind_max: entero)
    |\br{cpi}|izquierda = |\textrm{copia de números desde}| ind_min |\textrm{hasta}| ind_mitad
    derecha = |\textrm{copia de números desde}| ind_mitad+1 |\textrm{hasta}| ind_max|\br{cpf}|
    i = 0
    j = 0

    |\br{ai}|while i < izquierda.longitud and j < derecha.longitud
        if izquierda[i] |\bxl{es}|<=|\bxr{es}| derecha[j]
            números[ind_min+i+j] = izquierda[i]
            i++
        else
            números[ind_min+i+j] = derecha[j]
            j++|\br{af}|

    |\br{ii}|while i < izquierda.longitud
        números[ind_min+i+j] = izquierda[i]
        i++|\br{if}|

    |\br{di}|while j < derecha.longitud
        números[ind_min+i+j] = derecha[j]
        j++|\br{df}|
\end{pseudocode}
\begin{annotations}
    \codebrace[width=5.5cm]{cpi}{cpf}{Se crean arreglos nuevos con cada mitad. La complejidad espacial es entonces $\Theta(n)$.}
    \codebrace[xshift=7.8cm, width=4.5cm]{ai}{af}{Lógica principal, cuando ambos subarreglos tienen elementos: se toma el menor y se añade a \texttt{números}.}
    \codebrace[xshift=6cm, width=7cm]{ii}{if}{Si quedan elementos en el arreglo de la izquierda, ya no quedan en el de la derecha. Se añade el excedente de la izquierda a \texttt{números}.}
    \codebrace[xshift=6cm, width=7cm]{di}{df}{Se añade el excedente de la derecha a \texttt{números}. Si este while corre, es porque el anterior no corrió y viceversa.}
    \codebox[notepos={(cbox@es.east)+(2cm,0.8cm)}, curve={bend left=25}, width=5cm]{es}{\texttt{<=} en lugar de \texttt{<} para estabilidad.}
\end{annotations}

Un algoritmo es \emph{estable} si mantiene el orden relativo de los elementos iguales. O sea, en caso de empate, deja primero al que estaba primero en el arreglo original. En la línea del comentario sobre estabilidad, se elige el elemento de la izquierda en caso de empate para que continúe ubicado antes que el de la derecha en la solución final.

% TODO: Análisis de complejidad? Así sea breve
% Ojo: para poder decir que es logarítmico formalmente, necesito el teorema maestro... o no?

\practicar[leetcode]{Sort an Array}{https://leetcode.com/problems/sort-an-array/}